\chapter{Angular identifiability of spherical harmonics}
\label{ch:method2}

\begin{figure}[htbp]
  \centering
  \includegraphics[width=0.99\linewidth]{figures/d7-sh-coverage.pdf}
  \caption{The argument of this chapter, before the algebra. A primitive's
  colour is allowed to depend on the direction it is viewed from, and that
  dependence is stored as a short series in the viewing direction. A
  degree-$\ell$ term oscillates $\ell$ times around the sphere, so telling the
  degrees apart requires seeing the primitive from enough different directions.
  Seen from all around (left) that is easy. Seen from an eleven-degree window
  (right) it is impossible: rescaled to the window, all three curves are the
  same gentle arc, and no fit can say which of them the data was asking for.
  Fitting all three anyway does not produce an error; it produces three
  coefficients that trade off against one another arbitrarily.}
  \label{fig:shcoverage}
\end{figure}

The same argument applies on the sphere of directions. Each primitive carries
view-dependent colour as spherical harmonics to degree $3$: $16$ coefficients per
channel, $48$ floats, against $11$ for position, scale, rotation and opacity
combined. Non-DC spherical harmonics are $45$ of $59$ floats --- $76\,\%$ of the
model. They are fitted from however many directions happen to see the primitive,
which may be very few.

Let $y(d)$ be the vector of SH basis functions evaluated at direction $d$. The
angular Gram matrix accumulated over views is
\begin{equation}
  A_k \;=\; \sum_n w_{nk}\, y(d_{nk})\, y(d_{nk})^{\!\top}.
  \label{eq:gram}
\end{equation}
A coefficient block is determinable only where the corresponding block of $A_k$
is well conditioned. The criterion is to retain degree $\ell$ for primitive $k$
only where
\begin{equation}
  \ell_{\max}(k) \;=\; \max\big\{\, \ell \;:\; \lambda_{\min}\big(A_k(\ell)\big) > \tau \,\big\},
  \label{eq:lmax}
\end{equation}
and to mask higher degrees to zero. A primitive seen from three directions
cannot support degree $3$; fitting it produces coefficients that encode noise
and, worse, produce visible artefacts under novel views that fall outside the
observed cone.

\keybox{why this is a separate contribution}{%
The angular mask is independent of the anisotropic filter. It survives even if the floor-occupancy diagnostic of
\S\ref{sec:instruments} kills the positional band-limit: the two share a principle --- fit only what the capture determines --- and they
share neither a mechanism nor a failure mode. The angular mask also opens a memory result the anisotropic filter has no route
to, since masked coefficients can go unstored.}

\section{Normalising the threshold}
\label{sec:b2scale}

Equation~\eqref{eq:lmax} as written thresholds $\lambda_{\min}(A_k(\ell))$
against an absolute $\tau$, and that does not survive contact with a real
capture. Two separate scale dependencies are at fault, and they need separate
fixes.

\paragraph{View count.} $A_k$ is a \emph{sum} over views, so a primitive seen by
$100$ cameras has eigenvalues an order of magnitude larger than one seen by $10$,
before any question of conditioning arises. The normalised form
\begin{equation}
  \bar A_k \;=\; \frac{A_k}{\sum_n w_{nk}}
  \label{eq:gramnorm}
\end{equation}
removes it: $\bar A_k$ is a weighted \emph{average} of $y y^{\!\top}$ over the
directions that saw the primitive, so $\tau$ means the same thing for primitives
with different view counts. Without this, $\tau$ silently encodes a preference
for well-observed primitives that has nothing to do with whether their angular
sampling is well spread.

\paragraph{Angular spread.} Normalising is necessary, and on its own it is insufficient. Even on
$\bar A_k$, $\lambda_{\min}$ still scales with how widely the directions are
spread, so one $\tau$ cannot serve a $100$-camera orbit and a $10$-camera cone.
Section~\ref{sec:b2amend} measures this: at $\tau = 0.01$ the absolute form keeps
everything on a full orbit and nothing on either degraded protocol, and at
$\tau = 0.001$ it keeps $82\,\%$ on the arc. It is a cliff. The scale-free condition number
\begin{equation}
  \ell_{\max}(k) \;=\; \max\Big\{\, \ell \;:\;
  \frac{\lambda_{\min}\big(\bar A_k(\ell)\big)}{\lambda_{\max}\big(\bar A_k(\ell)\big)} > \tau \,\Big\}
  \label{eq:lmaxcond}
\end{equation}
behaves as the criterion should at a single threshold, and is what the
implementation uses. Equation~\eqref{eq:lmax} is retained above because it is
the form the argument is most naturally stated in, and because the amendment is
a result: it is a change to the method that only measurement could have
produced, and \S\ref{sec:b2amend} reports the measurement that forced it, with the
derivation amended in the open.

\paragraph{Status: implemented, measured, and refuted.} The angular mask is evaluated on all
eight Blender scenes and across the six capture protocols of
\S\ref{sec:stress} without a GPU, since Equation~\eqref{eq:gram} depends only on
primitive positions and training view directions; Table~\ref{tab:b2} adds the
rendered-quality and model-size columns.

The criterion does what this chapter says it does in the \emph{geometric}
sense: it retains everything on a full orbit, nothing on a pencil, and responds
gradually in between. Choosing \emph{which} coefficients to keep at a fixed budget is a separate
matter, and it is the claim the standard set below tests. Section~\ref{sec:btwo} reports the test and the criterion
fails it: against magnitude pruning restricted to the angular mask's own nested degrees, so that the structure and the budget are
the same and only the criterion differs, the angular mask loses by up to
$\btwoBlockedGapMax$\,dB. The derivation that follows is therefore presented as
a proposal that was carried through to measurement and did not survive it. It is
kept in full because the argument is sound and the conclusion is still
wrong, and the gap between the two is what a reader can learn from.

One caveat carries through to every statement about memory: masking coefficients
to zero does not shrink a fixed-width PLY, so \emph{no model-size saving is
claimed from these numbers}. Realising it needs variable-length SH storage,
which Chapter~\ref{ch:plan} schedules for Spring, and which, given
\S\ref{sec:btwo}, should be scheduled for a criterion that works.